\section{Lower Bound}
We aim to construct states $\ket{\varphi_C}$ and $\rho_C$ such that their measurement outcomes are close, while the fidelity is small enough so that the algorithm will reject.

\subsection{SIC-POVM states}
A POVM is said to be \textit{informationally complete} if the quantum state can be uniquely determined from the measurement outcome probabilities.
Since the real vector space of Hermitian matrices on $\mathbb{C}^d$ has real dimension $d^2$, any such POVM must have at least $d^2$ elements.
If we further require that the pairwise Hilbert--Schmidt inner products between distinct POVM elements are equal, we call such a POVM a SIC-POVM.
Equivalently, a SIC-POVM in $\mathbb{C}^d$ can be specified by $d^2$ pure states $\ket{\chi_{k}}$ such that 
\begin{align}\label{eq_202606101020}
\left|\bra{\chi_{i}}\ket{\chi_j}\right|^2=\frac{1}{d+1},\quad i\neq j.
\end{align}
The corresponding POVM element is
\begin{align}\label{eq_202606101437}
E_k=\frac{1}{d}\ket{\chi_k}\bra{\chi_k}.
\end{align}
\begin{remark}
It is still unknown whether SIC-POVMs exist in every dimension $d$~\cite{10.1063/1.1737053}.
\end{remark}
Luckily, in the qubit case $d=2$, we can construct a SIC-POVM as follows.
Let
\begin{align*}
r_1&=\frac{1}{\sqrt{3}}(1,1,1),\quad r_2=\frac{1}{\sqrt{3}}(1,-1,-1),\quad r_3=\frac{1}{\sqrt{3}}(-1,1,-1),\quad r_4=\frac{1}{\sqrt{3}}(-1,-1,1).
\end{align*}
Define the rank-one projection
\begin{align}\label{eq_202606101024}
\ket{\chi_k}\bra{\chi_k}=\frac{1}{2}\left(I_2+r_k\cdot\sigma\right)
\end{align}
where $\sigma$ is the Pauli vector.
Then, it is a valid POVM, since
\begin{align*}
\sum_{k=1}^4 E_k=I_2+\frac{1}{4}\left(\sum_{k=1}^4r_k\right)\cdot\sigma=I_2.
\end{align*}

To check \eqref{eq_202606101020}, we compute
\begin{align}\label{eq_202606101002}
\left|\bra{\chi_{i}}\ket{\chi_j}\right|^2&=\tr\left(\ket{\chi_i}\bra{\chi_i}\ket{\chi_j}\bra{\chi_j}\right)=\frac{1}{4}\tr\left(I_2+r_i\cdot\sigma+r_j\cdot\sigma+\left(r_i\cdot\sigma\right)\left(r_j\cdot\sigma\right)\right).
\end{align}
Note that 
\begin{align*}
\tr\left(r_i\cdot\sigma\right)=\sum_{\ell=1}^3r_{i\ell}\tr\sigma_{\ell}=0
\end{align*}
and
\begin{align*}
\tr\left(\left(r_i\cdot\sigma\right)\left(r_j\cdot\sigma\right)\right)&=\sum_{a,b=1}^3r_{ia}r_{jb}\tr\left(\sigma_a\sigma_b\right)=2\sum_{a,b=1}^3r_{ia}r_{jb}\delta_{ab}=2r_i\cdot r_j.
\end{align*}
Also, by computation,
\begin{align*}
r_i\cdot r_j=-\frac{1}{3},\quad i\neq j.
\end{align*}
Thus, \eqref{eq_202606101002} becomes
\begin{align*}
\left|\bra{\chi_{i}}\ket{\chi_j}\right|^2&=\frac{1}{2}\left(1+r_i\cdot r_j\right)=\frac{1}{3}=\frac{1}{2+1}.
\end{align*}
Hence, $\{E_k\}_{k=1}^4$ represents a valid SIC-POVM in dimension $2$.

We now show that, for any one-qubit measurement basis, at least one of the cross terms $\ket{\chi_1}\bra{\chi_b}$ is suppressed for some $b\in\{2,3,4\}$.
\begin{lemma}
For any orthonormal basis $\{\ket{\varphi},\ket{\varphi^\perp}\}$ of $\mathbb{C}^2$, there exists $b\in\{2,3,4\}$ such that
\begin{align}\label{eq_202606101458}
\left|\bra{\varphi}\ket{\chi_1}\bra{\chi_b}\ket{\varphi}\right|+\left|\bra{\varphi^\perp}\ket{\chi_1}\bra{\chi_b}\ket{\varphi^\perp}\right|\le\sqrt{\frac{8}{9}}.
\end{align}
\end{lemma}
\begin{proof}
Let the Bloch vector of $\ket{\varphi}$ be $v\in\mathbb{R}^3$.
Then, by orthogonality, the Bloch vector of $\ket{\varphi^\perp}$ is $-v$.
As a result, we have
\begin{align*}
\left|\bra{\varphi}\ket{\chi_k}\right|^2&=\tr\left(\ket{\varphi}\bra{\varphi}\ket{\chi_k}\bra{\chi_k}\right)=\frac{1}{4}\tr\left(I_2+v\cdot\sigma+r_k\cdot\sigma+\left(v\cdot\sigma\right)\left(r_k\cdot\sigma\right)\right)=\frac{1}{2}\left(1+v\cdot r_k\right)
\end{align*}
and
\begin{align*}
\left|\bra{\varphi^\perp}\ket{\chi_k}\right|^2&=\tr\left(\ket{\varphi^\perp}\bra{\varphi^\perp}\ket{\chi_k}\bra{\chi_k}\right)=\frac{1}{4}\tr\left(I_2-v\cdot\sigma+r_k\cdot\sigma-\left(v\cdot\sigma\right)\left(r_k\cdot\sigma\right)\right)=\frac{1}{2}\left(1-v\cdot r_k\right)
\end{align*}
for any $k$.
Let $\alpha=v\cdot r_1$ and $\beta=v\cdot r_k$.
Then, \eqref{eq_202606101458} is equivalent to
\begin{align*}
\left(\frac{1}{2}\sqrt{(1+\alpha)(1+\beta)}+\frac{1}{2}\sqrt{(1-\alpha)(1-\beta)}\right)^2=\frac{1}{2}\left(1+\alpha\beta+\sqrt{\left(1-\alpha^2\right)\left(1-\beta^2\right)}\right)\le\frac{8}{9}.
\end{align*}

By direct computations,
\begin{align*}
v\cdot\left(r_1-r_2\right)&=\frac{2}{\sqrt{3}}\left(v_2+v_3\right)\\
v\cdot\left(r_1-r_3\right)&=\frac{2}{\sqrt{3}}\left(v_1+v_3\right)\\
v\cdot\left(r_1-r_4\right)&=\frac{2}{\sqrt{3}}\left(v_1+v_2\right).
\end{align*}
Also,
\begin{align*}
\left(v_1+v_2\right)^2+\left(v_2+v_3\right)^2+\left(v_1+v_3\right)^2=1+\left(v_1+v_2+v_3\right)^2\ge1.
\end{align*}
Thus, there exists $b\in\{2,3,4\}$ such that
\begin{align*}
|\alpha-\beta|=\left|v\cdot\left(r_1-r_b\right)\right|\ge\frac{2}{\sqrt{3}}\cdot\frac{1}{\sqrt{3}}=\frac{2}{3}.
\end{align*}
Observe that
\begin{align*}
\sqrt{\left(1-\alpha^2\right)\left(1-\beta^2\right)}\le\frac{2-\alpha^2-\beta^2}{2}.
\end{align*}
Hence,
\begin{align*}
\frac{1}{2}\left(1+\alpha\beta+\sqrt{\left(1-\alpha^2\right)\left(1-\beta^2\right)}\right)\le\frac{1}{2}\left(2-\frac{1}{2}\left(\alpha-\beta\right)^2\right)\le\frac{1}{2}\left(2-\frac{1}{2}\cdot\frac{4}{9}\right)=\frac{8}{9}.
\end{align*}
This proves the lemma.
\end{proof}

\iffalse
\subsection{Proof of Theorem 2}
Let $\left\{\ket{\chi_i}\right\}_{i=0}^3$ be the qubit SIC-POVM states.
Let $N=\lfloor 2^{10^{-10}n}\rfloor$ and let $C=\left(c^1,\cdots,c^N\right)$ with $c^i\in\{0,1,2,3\}^n$.
Define
\[
	v_C=\frac{1}{\sqrt{N}}\sum_{t=1}^N\ket{\varphi_{c^t}}\quad\text{and}\quad\rho_C=\frac{1}{N}\sum_{t=1}^N\ket{\varphi_{c^t}}\bra{\varphi_{c^t}}
\]
where
\[
	\ket{\varphi_{c^t}}=\ket{\chi_{c^t_1}}\otimes\cdots\otimes\ket{\chi_{c^t_n}}.
\]
Let $\ket{\varphi_C}$ be the state in the direction of $v_C$.
Denote $\sigma_1=\ket{\varphi_C}\bra{\varphi_C}$ and $\sigma_2=\rho_C$.
By Lemmas 20 and 21, we have
\[
	d_{\mathrm{TV}}\left(M(\sigma_1),M(\sigma_2)\right)\le\frac{5}{N}+2^{-0.0001n}=2^{-\Omega (n)}
\]
for any $10^{-8}n$-adaptive basis $M$.

Consider any algorithm using $T=2^{o(n)}$ copies.
Let $\mathfrak{D}_1$ and $\mathfrak{D}_2$ be the distributions corresponding to measuring $2^T$ copies of $\ket{\varphi_C}\bra{\varphi_C}$ and $\rho_C$.
Denote the outcome $Y_{<t}=(Y_1,\cdots,Y_{t-1})$.
Therefore,
\begin{align*}
d_{\mathrm{TV}}\left(\mathfrak{D}_1,\mathfrak{D}_2\right)&\le\sum_{t=1}^Td_{\mathrm{TV}}\left(M_t(y_{<t})(\sigma_1),M_t(y_{<t})(\sigma_2)\right)\\
&\le\sum_{t=1}^T\sup_{y_{<t}}\ d_{\mathrm{TV}}\left(M_t(y_{<t})(\sigma_1),M_t(y_{<t})(\sigma_2)\right)\\
&\le T\cdot2^{-\Omega(n)}=2^{-\Omega(n)}.
\end{align*}
We conclude the whole proof.
\fi

%To bound the fidelity, we compute
%\begin{align*}
%\bra{\varphi_C}\rho_C\ket{\varphi_C}&=\frac{1}{N^2}%\sum\bra{\varphi_{c^s}}\ket{\varphi_{c^t}}%%\bra{\varphi_{c^t}}\ket{\varphi_{c^q}}
%\end{align*}
